\section{A sharper full-data residual bound}\label{h:end:tail-section}

The reflection estimate suffices for the main theorem. The following
bound gives a sharper success probability by identifying the full
conditional residual sphere. It allows arbitrary data-based Euclidean
outputs, including randomized ones.

\begin{lemma}[A memory-independent constant-accuracy bound]
\label{h:end:linear}
For all sufficiently large $d$, let $0\le T\le d/4$ be an integer.
Let $S$ be uniform on $S^{d-1}$, let $X$ be an independent
$T$-by-$d$ standard Gaussian matrix, and put $Y=XS$. For every
$0<\epsilon\le1/10$, any possibly randomized estimator $W\in\mathbb R^d$
based on the entire pair $(X,Y)$ satisfies
\begin{equation}\label{h:end:full-data-bound}
 \mathbb P\{\|W-S\|\le\epsilon\}
 \le \frac13+\left(\frac25\right)^{3d/4-1}
 <\frac23.
\end{equation}
Consequently, a learner with uniform-sphere angular success at least
$2/3$ at accuracy $\epsilon\le1/10$ must have $T>d/4$, even with
unlimited memory and access to all pre-generated samples.
\end{lemma}

\begin{proof}
Let $U$ be the row span of $X$, write $r=\dim U\le T$, and put
$k=d-r$. The labels determine $u=P_US$: the restriction of $X$ to
$U=(\ker X)^\perp$ is injective. More explicitly, when $T>0$,
diagonalizing $XX^{\mathsf T}$ shows that the consistent label $Y$
determines
\[
 u=\lim_{\eta\downarrow0}
 X^{\mathsf T}(XX^{\mathsf T}+\eta I_T)^{-1}Y.
\]
This also shows that $u$ is a measurable function of the data.
The analogous limit
\[
 P_U=\lim_{\eta\downarrow0}X^{\mathsf T}
 (XX^{\mathsf T}+\eta I_T)^{-1}X
\]
shows that the projection is measurable in $X$.
For the empty horizon $T=0$, $X$ has no rows and
$U=\{0\}$, $u=0$. The same values hold whenever the row span is
zero, including an all-zero matrix. In all these cases the residual
radius $\rho=\sqrt{1-\|u\|^2}$ equals one.

Standard Borel disintegration supplies a conditional probability
kernel for the actual joint law and the Borel data map
$(X,S)\mapsto(X,XS)$ \cite[Theorem~2.3]{GagnePanangaden2023}.
To identify it almost everywhere, condition first on any fixed matrix
$X$ with $r<d$, and generate $S$ as $G/\|G\|$ from an independent standard
Gaussian vector. Its orthogonal components $G_U$ and $G_{U^\perp}$
are independent. In polar coordinates on $U^\perp$, the Gaussian
volume element factors as a constant times
$e^{-\varrho^2/2}\varrho^{k-1}\,d\varrho\,d\omega$ for
$\varrho>0$, where $d\omega$ is spherical area. Hence
\[
 V=\frac{G_{U^\perp}}{\|G_{U^\perp}\|}
\]
is uniform on the unit sphere of $U^\perp$ and is independent of
$(G_U,\|G_{U^\perp}\|)$. Here $\|G_{U^\perp}\|>0$ almost surely
because $k=d-r\ge d-T>0$. The decomposition
\[
 S=u+\rho V,\qquad
 u=\frac{G_U}{\sqrt{\|G_U\|^2+\|G_{U^\perp}\|^2}},
 \qquad
 \rho=\frac{\|G_{U^\perp}\|}
 {\sqrt{\|G_U\|^2+\|G_{U^\perp}\|^2}}
\]
shows that $V$ is independent of $(u,\rho)$ conditional on $X$.
Since $u$ and $\rho$ are determined by $(X,Y)$, the conditional
law of the residual signal given the data is uniform on the sphere
of radius $\rho$ in $U^\perp$, for almost every data pair. When
$U=\{0\}$ this is the whole unit sphere and $\rho=1$. This argument
uses the Gaussian direction directly, including that zero-rank
case.

Independence of $X$ and $S$, together with
$\mathbb E SS^{\mathsf T}=I_d/d$, gives
\[
 \mathbb E\|u\|^2
 =\mathbb E\frac{\dim U}{d}
 \le\frac Td\le\frac14.
\]
Markov's inequality therefore yields
\begin{equation}\label{h:end:radius}
 \mathbb P\{\rho<1/2\}
 =\mathbb P\{\|u\|^2>3/4\}\le\frac13.
\end{equation}
On data with $\rho\ge1/2$, fix any proposed estimate
$w\in\mathbb R^d$. If $\|w-S\|\le\epsilon$, orthogonal projection
onto $U^\perp$ and division by $\rho$ imply
\[
 \left\|V-\frac{P_{U^\perp}w}{\rho}\right\|
 \le\frac{\epsilon}{\rho}\le\frac15.
\]
If this ball meets the unit sphere of $U^\perp$, choose a point of
the intersection. The entire intersection lies in the ball of
radius $2/5$ about that sphere point. Lemma~\ref{h:local:caps},
applied in ambient dimension $k\ge d-T\ge3d/4$, bounds its
conditional spherical probability by
\[
 \left(\frac25\right)^{k-1}
 \le\left(\frac25\right)^{3d/4-1}.
\]
An empty intersection has probability zero. This bound holds for
each $w$, so it also holds after averaging any output probability
kernel based on the data. Combining it with
\eqref{h:end:radius} proves \eqref{h:end:full-data-bound} for large
$d$. Finally, for unit vectors angular error at most $\epsilon$
implies Euclidean error at most $\epsilon$, since
$2\sin(\theta/2)\le\theta$. Granting all pre-generated rows lets a
full-data estimator simulate any earlier stopping rule, which proves
the stated consequence.
\end{proof}
